% Chapter 3: Higher-Order Linear Differential Equations
% Template file - Replace this content with your chapter material

%% Preamble - Common setup for all beamer presentations
%% Include this file in your chapter presentations

\documentclass[aspectratio=169, 12pt]{beamer}

% Packages
\usepackage[utf8]{inputenc}
\usepackage[T1]{fontenc}
\usepackage{amsmath}
\usepackage{amssymb}
\usepackage{amsthm}
\usepackage{mathtools}
\usepackage{graphicx}
\usepackage{xcolor}
\usepackage{tcolorbox}
\usepackage{tikz}
\usetikzlibrary{shadows, arrows, shapes}
\usepackage{hyperref}
\usepackage{listings}
\usepackage{geometry}

% Custom colors (import from custom_colors.tex if available)
\definecolor{defcolor}{RGB}{230, 245, 250}
\definecolor{defborder}{RGB}{0, 102, 204}
\definecolor{thmcolor}{RGB}{245, 230, 250}
\definecolor{thmborder}{RGB}{153, 51, 153}
\definecolor{excolor}{RGB}{240, 255, 240}
\definecolor{exborder}{RGB}{34, 139, 34}
\definecolor{probcolor}{RGB}{255, 250, 240}
\definecolor{probborder}{RGB}{255, 140, 0}

% Beamer Theme
\usetheme{Madrid}
\usecolortheme{default}
\setbeamertemplate{navigation symbols}{}
\setbeamertemplate{footline}[frame number]

% tcolorbox environment definitions

%% Definition Box
\newtcolorbox{definitionbox}[1]{
  colback=defcolor,
  colframe=defborder,
  coltitle=white,
  colbacktitle=defborder,
  title=#1,
  fonttitle=\bfseries\large,
  boxrule=2pt,
  arc=5pt,
  left=8pt,
  right=8pt,
  top=8pt,
  bottom=8pt,
  shadow={2pt}{-2pt}{0pt}{black!20},
  enhanced
}

%% Theorem Box
\newtcolorbox{theorembox}[1]{
  colback=thmcolor,
  colframe=thmborder,
  coltitle=white,
  colbacktitle=thmborder,
  title=#1,
  fonttitle=\bfseries\large,
  boxrule=2pt,
  arc=5pt,
  left=8pt,
  right=8pt,
  top=8pt,
  bottom=8pt,
  shadow={2pt}{-2pt}{0pt}{black!20},
  enhanced
}

%% Example Box
\newtcolorbox{examplebox}[1]{
  colback=excolor,
  colframe=exborder,
  coltitle=white,
  colbacktitle=exborder,
  title=#1,
  fonttitle=\bfseries\large,
  boxrule=2pt,
  arc=5pt,
  left=8pt,
  right=8pt,
  top=8pt,
  bottom=8pt,
  shadow={2pt}{-2pt}{0pt}{black!20},
  enhanced
}

%% Problem Box
\newtcolorbox{problembox}[1]{
  colback=probcolor,
  colframe=probborder,
  coltitle=white,
  colbacktitle=probborder,
  title=#1,
  fonttitle=\bfseries\large,
  boxrule=2pt,
  arc=5pt,
  left=8pt,
  right=8pt,
  top=8pt,
  bottom=8pt,
  shadow={2pt}{-2pt}{0pt}{black!20},
  enhanced
}

% Custom commands
\newcommand{\RR}{\mathbb{R}}
\newcommand{\CC}{\mathbb{C}}
\newcommand{\NN}{\mathbb{N}}
\newcommand{\ZZ}{\mathbb{Z}}
\newcommand{\dd}[1]{\frac{d}{d#1}}
\newcommand{\dydx}{\frac{dy}{dx}}
\newcommand{\dd2ydx2}{\frac{d^2y}{dx^2}}

% Title page formatting
\title[Your Chapter Title]{Your Chapter Title}
\author{Author Name}
\institute{Institution}
\date{\today}

% Remove the footer from title page
\setbeamertemplate{footline}{}


\title[Chapter 3]{Chapter 3: Higher-Order Linear Differential Equations}
\subtitle{Methods and Solutions}
\author{Your Name}
\date{\today}

\begin{document}

\begin{frame}
    \titlepage
\end{frame}

\begin{frame}
    \frametitle{Chapter Overview}
    \tableofcontents
\end{frame}

\section{Linear Equations}

\begin{frame}
    \frametitle{General Theory}
    
    \begin{definitionbox}{Linear Homogeneous Equation}
        An $n$-th order linear homogeneous differential equation has the form:
        \[
        a_n(x)\ddnydxn{n} + a_{n-1}(x)\ddnydxn{n-1} + \cdots + a_1(x)\dydx + a_0(x)y = 0
        \]
    \end{definitionbox}
\end{frame}

\section{Constant Coefficients}

\begin{frame}
    \frametitle{Second-Order with Constant Coefficients}
    
    \begin{definitionbox}{Standard Form}
        The equation
        \[
        a\dd2ydx2 + b\dydx + cy = 0
        \]
        has the \textbf{characteristic equation}:
        \[
        ar^2 + br + c = 0
        \]
    \end{definitionbox}
\end{frame}

\begin{frame}
    \frametitle{Solution Cases}
    
    \begin{theorembox}{Case 1: Distinct Real Roots}
        If $r_1 \neq r_2$ are real, then the general solution is:
        \[
        y = c_1 e^{r_1 x} + c_2 e^{r_2 x}
        \]
    \end{theorembox}
    
    \pause
    
    \begin{theorembox}{Case 2: Repeated Real Root}
        If $r_1 = r_2 = r$, then:
        \[
        y = c_1 e^{rx} + c_2 x e^{rx}
        \]
    \end{theorembox}
\end{frame}

\begin{frame}
    \frametitle{Solution Cases (continued)}
    
    \begin{theorembox}{Case 3: Complex Conjugate Roots}
        If $r = \alpha \pm i\beta$, then:
        \[
        y = e^{\alpha x}(c_1 \cos \beta x + c_2 \sin \beta x)
        \]
    \end{theorembox}
\end{frame}

\section{Examples}

\begin{frame}
    \frametitle{Example 1: Distinct Real Roots}
    
    \begin{problembox}{Problem}
        Solve: $\dd2ydx2 - 3\dydx + 2y = 0$ with $y(0) = 1$, $y'(0) = 0$
    \end{problembox}
\end{frame}

\begin{frame}
    \frametitle{Solution}
    
    \begin{examplebox}{Solution}
        Characteristic equation: $r^2 - 3r + 2 = 0$
        
        Factoring: $(r-1)(r-2) = 0$, so $r_1 = 1, r_2 = 2$
        
        General solution: $y = c_1 e^x + c_2 e^{2x}$
        
        Using $y(0) = 1$: $c_1 + c_2 = 1$
        
        Using $y'(0) = 0$: $c_1 + 2c_2 = 0$
        
        Solving: $c_1 = 2, c_2 = -1$
        
        \textbf{Particular solution:} $y = 2e^x - e^{2x}$
    \end{examplebox}
\end{frame}

\begin{frame}
    \frametitle{Example 2: Complex Roots}
    
    \begin{problembox}{Problem}
        Solve: $\dd2ydx2 + 2\dydx + 2y = 0$
    \end{problembox}
\end{frame}

\begin{frame}
    \frametitle{Solution}
    
    \begin{examplebox}{Solution}
        Characteristic equation: $r^2 + 2r + 2 = 0$
        
        Using the quadratic formula:
        \[
        r = \frac{-2 \pm \sqrt{4-8}}{2} = \frac{-2 \pm 2i}{2} = -1 \pm i
        \]
        
        So $\alpha = -1, \beta = 1$
        
        General solution:
        \[
        y = e^{-x}(c_1 \cos x + c_2 \sin x)
        \]
    \end{examplebox}
\end{frame}

\section{Summary}

\begin{frame}
    \frametitle{Summary}
    
    \begin{itemize}
        \item Higher-order linear equations with constant coefficients are solved via characteristic equation
        \item Solutions depend on whether roots are distinct, repeated, or complex
        \item Initial conditions determine the constants
    \end{itemize}
\end{frame}

\end{document}
